%! Parametrization of Implicitly Defined Functions — Unit 4, Lesson 4.7 — Exit Ticket
\documentclass[10pt]{article}
\usepackage{precalculus-article}
\usepackage{precalculus-boxes}
\newcommand{\TallMath}[1]{$\displaystyle #1\rule[-1.4em]{0pt}{3.2em}$}

\begin{document}

\pageheader{Unit 4, Lesson 4.7}{Exit Ticket}
\namedateperiod

\vspace{0.12in}

\begin{notesbox}{Parametrization Check}

The ellipse $\dfrac{(x+1)^2}{9} + \dfrac{(y-2)^2}{16} = 1$ is implicitly defined.

\vspace{10pt}
\textbf{1.}\quad Identify $h$, $k$, $a$, and $b$ from the equation. Then write the
parametric equations that trace the ellipse counterclockwise for $0 \le t \le 2\pi$.

\vspace{6pt}
$h = $ \blank{1.5cm} \quad $k = $ \blank{1.5cm} \quad $a = $ \blank{1.5cm} \quad $b = $ \blank{1.5cm}

\vspace{10pt}
$x(t) = $ \blank{5cm} \qquad $y(t) = $ \blank{5cm}

\vspace{12pt}
\textbf{2.}\quad Using your parametrization from Question 1, find the coordinates of the
topmost and bottommost points on the ellipse. Show how you use $y(t)$ to find them.

\writelines{4}

\end{notesbox}

\end{document}
